\appendix
\section{Layer-wise Architecture Tables}
\label{app:architectures}
The alias \texttt{f} means \texttt{features}; \texttt{head} means
\texttt{classifier}. Shapes include batch size one. Input/output shapes expose
all channel counts. K/S/P denotes kernel/stride/padding; symmetric pairs use
\texttt{x}. A dash denotes not applicable. All listed parameters are trainable.
NCHW means batch, channels, height and width; thus each row gives input and
output channel counts. A convolution has as many filters as output channels.
ReLU rows show activation placement; convolution and linear entries have no
implicit activation. BN denotes batch normalization
and GAP denotes adaptive global average pooling. Model B depthwise convolutions
use groups equal to input channels; its pointwise and all standard convolutions
use one group. DWConv and PWConv denote depthwise and pointwise convolutions.
Full names and groups are in the machine-readable architecture CSVs.
\subsection{Model A}
Table~\ref{tab:model_a_layers} gives the complete standard CNN, including
normalization scales/offsets and the classifier bias in its parameter sum.
\begingroup\scriptsize\setlength{\tabcolsep}{3pt}
\begin{longtable}{lllllr}
\caption{Model A leaf layers, shapes and trainable parameter counts.}\label{tab:model_a_layers}\\
\toprule Layer & Operation & Input (NCHW) & Output & K/S/P & Params \\ \midrule\endfirsthead
\multicolumn{6}{l}{Model A layers (continued)}\\\toprule Layer & Operation & Input (NCHW) & Output & K/S/P & Params \\ \midrule\endhead
f.0.0 & Conv2d & $1\times3\times64\times64$ & $1\times24\times64\times64$ & 3x3/1x1/1x1 & 648 \\
f.0.1 & BN & $1\times24\times64\times64$ & $1\times24\times64\times64$ & -/-/- & 48 \\
f.0.2 & ReLU & $1\times24\times64\times64$ & $1\times24\times64\times64$ & -/-/- & 0 \\
f.0.3 & MaxPool & $1\times24\times64\times64$ & $1\times24\times32\times32$ & 2/2/0 & 0 \\
f.1.0 & Conv2d & $1\times24\times32\times32$ & $1\times48\times32\times32$ & 3x3/1x1/1x1 & 10368 \\
f.1.1 & BN & $1\times48\times32\times32$ & $1\times48\times32\times32$ & -/-/- & 96 \\
f.1.2 & ReLU & $1\times48\times32\times32$ & $1\times48\times32\times32$ & -/-/- & 0 \\
f.1.3 & MaxPool & $1\times48\times32\times32$ & $1\times48\times16\times16$ & 2/2/0 & 0 \\
f.2.0 & Conv2d & $1\times48\times16\times16$ & $1\times96\times16\times16$ & 3x3/1x1/1x1 & 41472 \\
f.2.1 & BN & $1\times96\times16\times16$ & $1\times96\times16\times16$ & -/-/- & 192 \\
f.2.2 & ReLU & $1\times96\times16\times16$ & $1\times96\times16\times16$ & -/-/- & 0 \\
f.2.3 & MaxPool & $1\times96\times16\times16$ & $1\times96\times8\times8$ & 2/2/0 & 0 \\
f.3.0 & Conv2d & $1\times96\times8\times8$ & $1\times128\times8\times8$ & 3x3/1x1/1x1 & 110592 \\
f.3.1 & BN & $1\times128\times8\times8$ & $1\times128\times8\times8$ & -/-/- & 256 \\
f.3.2 & ReLU & $1\times128\times8\times8$ & $1\times128\times8\times8$ & -/-/- & 0 \\
f.3.3 & MaxPool & $1\times128\times8\times8$ & $1\times128\times4\times4$ & 2/2/0 & 0 \\
pool & GAP & $1\times128\times4\times4$ & $1\times128\times1\times1$ & -/-/- & 0 \\
flatten & Flatten & $1\times128\times1\times1$ & $1\times128$ & -/-/- & 0 \\
head & Linear & $1\times128$ & $1\times10$ & -/-/- & 1290 \\
\midrule\multicolumn{5}{r}{Total trainable parameters} & 164,962 \\ \bottomrule
\end{longtable}
\endgroup

\clearpage
\subsection{Model B}
Table~\ref{tab:model_b_layers} gives the complete depthwise-separable CNN.
The layer sum independently verifies the parameter budget.
\begingroup\scriptsize\setlength{\tabcolsep}{3pt}
\begin{longtable}{lllllr}
\caption{Model B leaf layers, shapes and trainable parameter counts.}\label{tab:model_b_layers}\\
\toprule Layer & Operation & Input (NCHW) & Output & K/S/P & Params \\ \midrule\endfirsthead
\multicolumn{6}{l}{Model B layers (continued)}\\\toprule Layer & Operation & Input (NCHW) & Output & K/S/P & Params \\ \midrule\endhead
f.0.0 & Conv2d & $1\times3\times64\times64$ & $1\times24\times64\times64$ & 3x3/1x1/1x1 & 648 \\
f.0.1 & BN & $1\times24\times64\times64$ & $1\times24\times64\times64$ & -/-/- & 48 \\
f.0.2 & ReLU & $1\times24\times64\times64$ & $1\times24\times64\times64$ & -/-/- & 0 \\
f.0.3 & MaxPool & $1\times24\times64\times64$ & $1\times24\times32\times32$ & 2/2/0 & 0 \\
f.1.0.0 & DWConv & $1\times24\times32\times32$ & $1\times24\times32\times32$ & 3x3/1x1/1x1 & 216 \\
f.1.0.1 & BN & $1\times24\times32\times32$ & $1\times24\times32\times32$ & -/-/- & 48 \\
f.1.0.2 & ReLU & $1\times24\times32\times32$ & $1\times24\times32\times32$ & -/-/- & 0 \\
f.1.0.3 & PWConv & $1\times24\times32\times32$ & $1\times48\times32\times32$ & 1x1/1x1/0x0 & 1152 \\
f.1.0.4 & BN & $1\times48\times32\times32$ & $1\times48\times32\times32$ & -/-/- & 96 \\
f.1.0.5 & ReLU & $1\times48\times32\times32$ & $1\times48\times32\times32$ & -/-/- & 0 \\
f.1.1 & MaxPool & $1\times48\times32\times32$ & $1\times48\times16\times16$ & 2/2/0 & 0 \\
f.2.0.0 & DWConv & $1\times48\times16\times16$ & $1\times48\times16\times16$ & 3x3/1x1/1x1 & 432 \\
f.2.0.1 & BN & $1\times48\times16\times16$ & $1\times48\times16\times16$ & -/-/- & 96 \\
f.2.0.2 & ReLU & $1\times48\times16\times16$ & $1\times48\times16\times16$ & -/-/- & 0 \\
f.2.0.3 & PWConv & $1\times48\times16\times16$ & $1\times96\times16\times16$ & 1x1/1x1/0x0 & 4608 \\
f.2.0.4 & BN & $1\times96\times16\times16$ & $1\times96\times16\times16$ & -/-/- & 192 \\
f.2.0.5 & ReLU & $1\times96\times16\times16$ & $1\times96\times16\times16$ & -/-/- & 0 \\
f.2.1 & MaxPool & $1\times96\times16\times16$ & $1\times96\times8\times8$ & 2/2/0 & 0 \\
f.3.0.0 & DWConv & $1\times96\times8\times8$ & $1\times96\times8\times8$ & 3x3/1x1/1x1 & 864 \\
f.3.0.1 & BN & $1\times96\times8\times8$ & $1\times96\times8\times8$ & -/-/- & 192 \\
f.3.0.2 & ReLU & $1\times96\times8\times8$ & $1\times96\times8\times8$ & -/-/- & 0 \\
f.3.0.3 & PWConv & $1\times96\times8\times8$ & $1\times128\times8\times8$ & 1x1/1x1/0x0 & 12288 \\
f.3.0.4 & BN & $1\times128\times8\times8$ & $1\times128\times8\times8$ & -/-/- & 256 \\
f.3.0.5 & ReLU & $1\times128\times8\times8$ & $1\times128\times8\times8$ & -/-/- & 0 \\
f.3.1 & MaxPool & $1\times128\times8\times8$ & $1\times128\times4\times4$ & 2/2/0 & 0 \\
pool & GAP & $1\times128\times4\times4$ & $1\times128\times1\times1$ & -/-/- & 0 \\
flatten & Flatten & $1\times128\times1\times1$ & $1\times128$ & -/-/- & 0 \\
head & Linear & $1\times128$ & $1\times10$ & -/-/- & 1290 \\
\midrule\multicolumn{5}{r}{Total trainable parameters} & 22,426 \\ \bottomrule
\end{longtable}
\endgroup

